\documentclass[11pt]{article}

\usepackage{fontspec}
\setmainfont{TeX Gyre Pagella}

\usepackage{amsmath,amssymb,amsthm}
\usepackage[margin=1.25in]{geometry}
\usepackage{hyperref}
\usepackage{listings}
\usepackage{xcolor}
\usepackage{longtable}
\usepackage{array}

\theoremstyle{definition}
\newtheorem{definition}{Definition}
\newtheorem{principle}{Design Principle}
\newtheorem{incident}{Incident}

\lstset{
  basicstyle=\ttfamily\small,
  breaklines=true,
  frame=single,
  columns=fullflexible
}

\hypersetup{
  colorlinks=true,
  linkcolor=black,
  urlcolor=black,
  pdftitle={Specification of the Distinction-First Pipeline},
  pdfauthor={Flyxion}
}

\title{Specification of the Distinction-First Pipeline: \\
Continuation, Accretion, and Externalized Memory \\
in a Local-Model Essay-Writing System}
\author{Flyxion \\ \small Independent Researcher}
\date{\today}

\begin{document}
\maketitle

\begin{abstract}
This essay specifies a nine-stage pipeline for producing chapter-length
LaTeX prose with a locally hosted language model, and gives the
rationale and full mechanics of every stage. The system is organized
around a single claim: its purpose is continuation, not generation.
Every stage exists to treat something already established --- a
distinction, a definition, an intent, an earlier chapter's conclusion
--- as a constraint on what a later stage may produce, rather than to
produce content from an unconstrained space. This follows the same
commitments that organize the author's broader theoretical program:
history before state, repair before correctness, admissibility before
prediction. It also follows from nine more specific, independently
observed habits in how the author actually writes: essays begin from a
distinction judged worth preserving rather than from a thesis; writing
functions as a discovery process rather than a report of something
already known; memory is externalized into structure rather than held
in mind; a corpus grows by accretion, one paper becoming a chapter
becoming a monograph becoming a sequence; there is a recurring search
for what survives a transformation rather than a fixed claim proved for
its own sake; recurring criticism is promoted into permanent mechanism
rather than treated as a local fix; concrete examples function as
witnesses of an already-established abstract structure rather than as
its motivation; long works grow by discovering a missing dependency and
writing it as the next chapter, rather than by planning a fixed table
of contents in advance; and underneath all of this sits a persistent
concern with how a large structure maintains its identity while
changing, as opposed to simply producing more of it. Nine stages
implement these commitments: a distinction-extraction stage that
replaces thesis-first outlining; an exploratory drafting stage in which
the main claim crystallizes only at the end; an admissibility gate that
tests for a locatable distinction rather than a locatable thesis; a
dependency-discovery stage that queues missing concepts as new chapters
instead of patching over them; a resonance-ranked, pin-overridden,
ledger-augmented corpus cross-referencing stage; a critique-revision
loop that distinguishes genuine convergence from mere stalling and
checks explicitly for a named invariant; a witness-finding stage that
constructs concrete instances of the draft's abstract structure only
after that structure is finished; and a fidelity audit against the
original distinction. A persistent ledger and a persistent critique log
externalize memory across runs, and a companion driver script accretes
a small graph of chapters by recursively processing whatever
dependencies a given chapter discovers it needs. Four of the nine
stages were derived by taking vocabulary from the Phoenix Protocol, a
wave-interference memory architecture under separate audit, and
stripping it of its physical claims, keeping only whatever functioned
as a legitimate control pattern for an iterative language-model loop.
One genuine implementation bug, found through testing, is documented as
an incident report rather than silently fixed, since the failure mode
it represents --- an identity mechanism that collapses under
composition --- is exactly the kind of thing this system exists to
catch in prose, and very nearly failed to catch in its own source.
\end{abstract}

\section{Motivation}

A large interconnected research program accumulates a vocabulary faster
than any single author can hold consistently in mind across hundreds of
pages written over months. Terms such as admissibility, repair
operator, witness, and continuation acquire precise meanings in early
essays and then risk being redefined, loosened, or silently repurposed
in later ones, not out of carelessness but simply because no single
essay-writing session has the whole corpus in view. A tool that assists
with drafting can either ignore this problem, in which case it
accelerates the production of prose while doing nothing about the risk
of terminological drift, or it can be built to actively check new
material against the existing corpus, against the author's own past
conclusions, and against the author's own stated intent, as part of its
ordinary operation. This pipeline takes the second approach throughout.

A second and more specific motivation concerns the actual shape of how
these essays get written, which was reconstructed by reading across the
author's essays, monographs, code repositories, and the way the author
has discussed writing over many conversations, rather than by reasoning
abstractly about what a good drafting tool should do. That reading
surfaced nine recurring habits, true of the author's writing long
before any of it was language-model-assisted, that a generic
draft-critique-revise loop does not capture on its own.

Essays do not begin from a thesis. They begin from a distinction judged
worth preserving --- a repair operation, a difference between storage
and trace, a concern about reconstruction versus recovery --- that
looks, on its own, too small to support a paper. The thesis is what the
essay finds by taking that distinction seriously at length, not what it
starts from. Writing itself functions as a discovery process rather
than a report of something already known: the essays read as ``I
suspect X, let's see what survives formalization'' rather than ``I know
X and will explain it,'' which is why they typically accumulate
definitions, propositions, and examples before arriving at anything
that would traditionally be called the main claim. Memory is repeatedly
externalized into structure rather than held in the author's head:
theorem environments, chapter hierarchies, named primitives, visible
draft histories, and large collections of essays rather than one
continuously growing document all serve this function. There is a
strong tendency toward corpus construction rather than essay
production: a single paper becomes a chapter, a chapter becomes a
monograph, a monograph becomes part of a sequence, and terminological
drift matters enormously for a structure like that in a way it would
not for an isolated essay. There is a recurring search for invariants
--- what survives reconstruction, projection, correction, compression,
memory loss, a change of coordinates --- in place of theorems proved
for their own sake. Criticism is repeatedly promoted into architecture
rather than treated as a local correction: a chapter that drifts from
its intent does not get fixed once, it produces a fidelity check that
runs on every future chapter; a revision loop that stalls does not get
manually interrupted once, it produces a general stall detector.
Examples function as witnesses of an already-established abstract
structure rather than as the motivating case that produced it: the
abstract machinery typically comes first, and significant effort goes
into finding or constructing examples that demonstrate it afterward.
Long works grow by accretion rather than advance planning: one chapter
reveals a missing dependency, the dependency becomes the next chapter,
which reveals another missing dependency behind it, so that a
three-hundred-page monograph looks less like a tree planned from the
root down and more like a graph that gradually reveals its own hidden
nodes. And underneath all of this sits a persistent distinction between
generation and continuation --- between producing new objects and
understanding how a large structure maintains its identity while
changing over time --- that shows up in the author's theoretical work,
software workflows, critique loops, treatment of memory, and now in
this pipeline's own architecture.

This pipeline is an attempt to build a system that matches that
description mechanically, stage by stage, rather than one that merely
drafts prose and critiques it.

\section{Three Levels of Description}
\label{sec:levels}

Before proceeding, it is worth separating three distinct levels at
which the claims in this essay operate, since later sections move
between them without always announcing the shift, and conflating them
would misstate what has actually been established.

The first level is the concrete shell implementation: a specific
script, a specific configuration file format, specific prompt templates
piped to a specific local model over standard input. Claims at this
level are falsifiable by reading the source or running it, and are
exactly as general as the code, no more.

The second level is a general architecture for language-model-assisted
writing: the pattern of separating drafting from critique, gating
cheaply before revising expensively, tracking a scalar signal to
distinguish convergence from stalling, discovering missing dependencies
rather than patching around them, and auditing a finished artifact
against original intent rather than only against intermediate critique.
This level is a claim about a reusable design, not about this
particular script; the shell implementation is one instance of it.

The third level is methodological: a claim about theoretical research
programs generally, and about this author's writing process
specifically, namely that a corpus accumulates vocabulary,
argumentative commitments, and structural dependencies that later
writing ought to be checked against, discover, and declare its
relationship to, rather than being free to override, ignore, or invent
silently. This level is the weakest to defend and the one this essay is
least entitled to assert broadly, since it is argued here from a single
case rather than established in general. Where the essay speaks in this
register it should be read as a hypothesis about this pipeline and
about programs and authors structured the way this one is, not as an
established claim about writing in general.

\section{Overview of the Nine-Stage Pipeline}
\label{sec:overview}

\begin{longtable}{@{}p{0.16\textwidth}p{0.78\textwidth}@{}}
\textbf{Stage} & \textbf{Purpose} \\ \hline
\endhead
1. Distinguish & Extract the seed distinction (A vs B) and the invariant
question it raises from raw, possibly messy notes. \\
2. Outline & Build a section skeleton from the distinction, with the
final section explicitly marked as an unresolved ``Crystallized Claim.'' \\
3. Draft & Write the chapter as an exploration: definitions and
propositions accumulate; the main claim is written last, in the
Crystallized Claim section, and may differ from anything anticipated
in the outline. \\
4. Ignite & Admissibility gate testing for a locatable distinction, not
a locatable thesis, plus definitions-before-use and populated sections. \\
5. Dependencies & Scan the draft for concepts used but not defined and
not already in the vocabulary; queue each as a new seed for a future
chapter rather than patching over the gap. \\
6. Rise & Resonance-ranked corpus cross-reference, with a pinning
override for under-repeated foundational terms, plus a persistent
ledger of every prior chapter's distinction. \\
7. Persist & Critique-revision loop tracking a severity score across
rounds, distinguishing convergence, dissolution (stalling), and
exhaustion (oscillation); critique checks explicitly for a named
invariant and logs its category tags to a cross-run log. \\
8. Witness & Find or construct concrete examples instantiating each
formal object in the finished draft, after the abstract structure is
settled, not before. \\
9. Audit & Compare the finished draft against the original seed
distinction for drift, then update the ledger and finalize the LaTeX
document. \\
\end{longtable}

A companion driver, discussed in Section~\ref{sec:accretion}, runs
stage 5's output back through the whole pipeline recursively, which is
what turns a single chapter-writing tool into something that accretes a
small graph of chapters from one seed.

\section{The Distinction, Not the Thesis}
\label{sec:distinguish}

The first stage exists to correct a specific mismatch between how a
generic outlining stage works and how these essays actually start. A
generic outlining stage asks a model to sketch an argument's structure,
which presupposes that an argument --- something with a thesis ---
already exists to be sketched. That is not how these essays begin. They
begin from something smaller and more specific: a distinction that on
its own looks too small to support a paper. The thesis is what the
essay finds by taking that distinction seriously at length, not what it
starts from, and an outlining stage that demands a thesis up front is
asking for something that, for this kind of writing, does not yet
exist at the point outlining happens.

The Distinguish stage is deliberately narrow. Given raw notes, it does
exactly one thing: name the distinction underneath them, as an
A-versus-B contrast, and name the invariant question that distinction
raises --- what survives some transformation, reconstruction,
projection, or correction, in a way that makes the distinction matter.
It is explicitly instructed not to summarize the notes and not to
produce anything resembling an outline; those are separate stages with
separate objectives, and folding them together would let the model
reach for a premature thesis while nominally only summarizing. The
output is two lines, machine-parseable, consumed by every subsequent
stage:

\begin{lstlisting}[caption={Distinguish stage output format}]
DISTINCTION: repair vs erasure
INVARIANT_QUESTION: What survives reconstruction from partial witnesses?
\end{lstlisting}

Everything downstream is built from this pair rather than from the raw
notes directly. The outline is built from the distinction with its
final section explicitly marked \texttt{TBD}. The draft is instructed
to explore the distinction rather than argue toward a conclusion. The
admissibility gate checks for the distinction's presence, not a
thesis's. The audit, at the end, checks fidelity to the distinction,
not to a thesis that was never supposed to exist yet at the point the
seed was written.

\section{Exploratory Drafting and the Crystallized Claim}
\label{sec:draft}

The draft prompt is built to match the discovery-over-reporting pattern
directly rather than to produce polished prose as quickly as possible.
It instructs the model to accumulate definitions and propositions that
follow from taking the seed distinction seriously, explicitly forbids
opening with a pre-announced conclusion such as ``this chapter argues
that,'' and requires a final subsection titled \emph{Crystallized
Claim} that states, in one paragraph and written last, whatever the
preceding exploration actually arrived at. This final section is
permitted, and expected, to be narrower than or different from anything
the outline anticipated, since the outline itself was built without a
thesis to outline toward in the first place.

\begin{principle}
A drafting prompt that asks a model to write toward a conclusion
already assumes the conclusion is known. A prompt whose only fixed
point is a distinction, with the conclusion deferred to a clearly
marked final section, mechanically reproduces writing as discovery
rather than writing as report, at the cost of the draft having no
guaranteed argumentative shape until the very last section is reached.
\end{principle}

This has a consequence for how every later stage treats the draft:
nothing before the Crystallized Claim section may be assumed to state
the chapter's actual position. This is why the admissibility gate in
Section~\ref{sec:gate} does not require a thesis in the opening, and
why the audit in Section~\ref{sec:audit} compares the finished draft to
the seed distinction rather than to any thesis stated partway through.

\section{The Admissibility Gate}
\label{sec:gate}

Before any revision effort is spent on a draft, the pipeline runs a
minimal gate against it. This gate is deliberately cheap, blunt, and
narrow, and its purpose is economy rather than quality control in the
usual sense. The critique-revision loop described in
Section~\ref{sec:persist} is comparatively expensive, running several
rounds of adversarial review against a document that may run to several
thousand words. Spending that effort on a draft that never had a
distinction to explore in the first place wastes it; no amount of
critique-and-revise cycling produces a chapter from an outline that was
never coherent to begin with, since the loop is designed to sharpen an
existing exploration, not to originate one. The gate is therefore
placed before the expensive stage and given the authority to reject
rather than merely warn.

What the gate checks is detectability, not quality, and this
distinction is worth stating precisely since it is easy to misread the
checks below as an evaluation of the draft's merit, which they are not.
The gate asks whether the system can locate a distinction-like object,
a term defined before use, and a populated section ending in a
Crystallized Claim --- not whether the distinction is interesting, the
definition is precise, or the argument in the populated section is
sound. A draft with a banal or ultimately unpromising distinction
passes this gate exactly as readily as a draft with an excellent one,
provided the distinction is stated plainly enough to be located. This
is intentional: judging the quality of an exploration is exactly the
kind of task the critique-revision loop exists to do at length, with
multiple rounds and an adversarial posture, and folding a quality
judgment into a single cheap gate call would either make the gate
expensive or make it unreliable.

Concretely, the gate checks four things. First, whether there is an
identifiable distinction being explored --- some contrast or pair the
draft is actually working with --- visible somewhere in the first
section, even where no conclusion has yet been drawn from it, since the
draft is explicitly permitted not to draw one there. Second, whether
every technical term used before it is formally introduced is in fact
defined somewhere in the draft, even informally. Third, whether the
draft has more than one section with actual argumentative content
rather than empty headers. Fourth, whether the draft in fact ends with
a Crystallized Claim section, since a draft that never reaches one has
failed to do the one thing the exploratory structure was for, no
matter how good its middle sections are. A draft that fails any of
these is regenerated, up to a small configurable number of retries,
before the pipeline gives up and asks for a better seed rather than
proceeding.

\begin{lstlisting}[caption={The admissibility check, stated as a predicate}]
def admissible(draft) -> bool:
    return (
        has_locatable_distinction(draft.first_section)   # not: has thesis
        and all(is_defined_before_use(t) for t in draft.terms)
        and count(sections_with_content(draft)) > 1
        and has_crystallized_claim_section(draft)
    )
# Not tested: whether the distinction is interesting, or the argument sound.
\end{lstlisting}

\section{Dependency Discovery and Chapter Accretion}
\label{sec:accretion}

This is the largest single mechanism in the pipeline, and it targets
the pattern of chapter accretion directly: long works in the author's
corpus are not planned as a fixed table of contents in advance, but
grow by writing one chapter, discovering that chapter depends on a
concept that has not itself been established, writing that concept as
another chapter, and repeating, so that a large monograph looks less
like a tree planned from the root down and more like a graph that
reveals its own missing nodes as it is written.

The Dependency stage implements the discovery half of this directly.
After a draft passes the admissibility gate, a separate model call
reads it specifically for unbuilt foundations rather than for prose
quality: technical terms used in a load-bearing way, not already in the
established vocabulary list, and not themselves defined anywhere in the
draft. Each such concept is a missing dependency --- something that
would need its own chapter before this one is fully grounded, even
though the current draft may otherwise be fine as exploratory writing.
Each one is reported with a one-sentence guess at what a chapter
establishing it would need to cover, and each is written out as a new
seed file rather than being flagged as a problem to fix in place:

\begin{lstlisting}[caption={A queued dependency seed file}]
CONTINUES: repair vs erasure
gyrofoil measure -- needs a chapter establishing its formal
definition and role in repair.
\end{lstlisting}

The \texttt{CONTINUES:} line matters and is discussed further in
Section~\ref{sec:continuation-decl}: the dependency is written down as
continuing the chapter that spawned it, not as a free-floating new
topic, preserving the provenance of why it was queued at all.

A companion script implements the accretion half. It runs the pipeline
on an initial seed, and whenever that run queues a dependency,
recursively runs the pipeline on that dependency as well, which may
itself queue further dependencies, until either the queue is empty or a
configured depth cap is reached. A processed-set guard prevents
infinite cycles if two chapters end up depending on each other, or the
same gap gets queued twice under different phrasing.

\begin{lstlisting}[caption={The accretion driver, simplified}]
def run_one(seed, depth):
    key = identity(seed)
    if key in processed: return       # cycle guard
    processed.add(key)
    if depth > MAX_DEPTH:
        return                        # leave queued, don't process
    pipeline_run(seed)                # may queue new dependencies
    for dep in queued_by(seed):
        run_one(dep, depth + 1)
\end{lstlisting}

\subsection{An Incident: Identity Collapsing Under Composition}
\label{sec:incident}

\begin{incident}
During testing of the accretion driver, a real bug surfaced that is
worth documenting rather than silently patching, since the failure
mode is a small instance of exactly the kind of thing this whole
system exists to catch --- an identity that fails to survive a
transformation.
\end{incident}

The pipeline's slug function, used to turn a chapter's title into a
filesystem-safe identifier, truncated its output to forty characters.
This was harmless for a single chapter, where the truncated slug only
ever needed to disambiguate a directory name already made unique by an
attached timestamp. It was not harmless for dependency filenames, which
are built by prefixing a parent chapter's slug onto a child's:
\texttt{parent-slug--child-term.seed}. Once a parent's own slug already
reached the forty-character cap, appending anything to it and
re-truncating discarded the appended part entirely, so a child's
filename collapsed onto its parent's. In practice this meant a
second-generation dependency was written to the exact path its parent
occupied, silently overwriting it, and the accretion driver --- which
tracks which files it still needs to move once processed --- attempted
to move a file that a nested recursive call had, unbeknownst to the
outer call, already moved a moment earlier, producing a file-not-found
error several stages removed from the actual cause.

The fix was not a patch at the point of failure but a second,
deliberately unbounded identity function used everywhere identity needs
to survive composition, while the truncated version was kept for the
one purpose it was always safe for:

\begin{lstlisting}[caption={Two identity functions with different survival guarantees}]
def slug(s):       # truncated -- safe only where a timestamp
    return dashify(s)[:40]   # already disambiguates, e.g. OUTDIR names

def slug_full(s):  # untruncated -- required wherever identity
    return dashify(s)        # must survive being prefixed onto again
\end{lstlisting}

The general lesson, stated at the level of Section~\ref{sec:levels}'s
third register: an identity mechanism that is safe in isolation is not
automatically safe under composition, and a system that recursively
builds identities out of other identities --- exactly what chapter
accretion does --- needs to ask, for every identity function it uses,
whether that function is closed under the composition the system
actually performs, not just correct for a single, non-recursive use.
This is the same question the pipeline asks of prose --- what survives
a transformation --- applied for once to its own source code rather
than to a chapter draft.

\section{Vocabulary as a Controlled Resource}
\label{sec:vocab}

The pipeline treats the project's established terminology as a
controlled vocabulary rather than as free text. A configuration file
lists every term that has an established meaning elsewhere in the
corpus --- named frameworks, together with core primitives such as
admissibility, distinction, witness, and continuation. This list is
injected verbatim into every prompt the pipeline constructs, with an
instruction to use these terms consistently and not to silently
redefine or repurpose them. The list is authored by hand and grows only
when a term becomes load-bearing somewhere in the corpus; the pipeline
does not infer or invent vocabulary on its own.

\begin{principle}
A system that assists with writing inside an established theoretical
program should be given the program's vocabulary as an explicit
constraint, not left to infer it from context on each invocation.
\end{principle}

This is a modest mechanism, but it changes what the critique stage in
Section~\ref{sec:persist} is capable of catching. Without an explicit
vocabulary, a critique pass can only evaluate a draft against itself:
internal contradictions, dropped threads, unsupported claims. With the
vocabulary supplied, the critique pass can additionally check a draft
against everything else the vocabulary has meant before, which is a
different and in some ways more important kind of error to catch,
since a single essay can be internally flawless while still
contradicting the definition of the same term two essays earlier.

Three claims are being made together here that are worth separating.
At the implementation level, the claim is only that a configuration
file lists terms and that this list is injected into prompts; this is
verifiable by reading the code. At the design level, the claim is a
recommendation --- that any system assisting with writing inside an
established body of work should be given that vocabulary explicitly
rather than left to infer it --- defensible on the general grounds that
inference from context is unreliable and an explicit list is cheap. At
the epistemic level, the claim is that meaning is actually preserved
across a corpus by this mechanism, which this essay has not
established and asserts only as a working hypothesis motivating the
design.

\section{Externalized Memory: Resonance, Pinning, and the Ledger}
\label{sec:corpus}

If a directory of existing essays is supplied, the pipeline performs a
cross-referencing pass before every critique. For every vocabulary term
that appears in both the configuration file and the current draft, the
system searches the existing corpus for other occurrences of that term
and counts them. Terms are then ranked by this count and only the
highest-ranked handful are actually pulled into the critique prompt,
together with a few lines of surrounding context from wherever they
occur in the corpus. This ranking step exists because an unranked
cross-reference pass does not scale: a corpus of any size produces far
more matching context than a language model's working context can
usefully absorb, and naively including all of it drowns the signal the
critique stage actually needs.

Ranking by occurrence count is a simple proxy for how load-bearing a
term is, and it is a proxy that costs nothing beyond a directory
search, rather than requiring any semantic understanding of the corpus
itself. It is also known to be wrong in one specific and important
direction: some foundational concepts are assumed throughout a corpus
rather than restated in it, and so appear rarely despite being the most
load-bearing terms present, while a term can be repeated often because
it is locally fashionable in a particular stretch of writing without
being central to the program as a whole. Frequency and importance need
not coincide, and a ranking step built on frequency alone would
silently favor the fashionable term over the assumed one.

This failure mode is addressed directly rather than merely noted. The
vocabulary configuration file supports a second category of entry, a
\emph{pinned} term, marked with a leading \texttt{!} in the source file
and stripped of that marker before the term is shown to the model.
Pinned terms bypass the occurrence-count ranking entirely: if a pinned
term appears in the current draft and has at least one hit anywhere in
the corpus, it is included in the critique context regardless of where
it would have fallen in the ranked list. Ranking-by-frequency thus
governs the unpinned majority of the vocabulary, where frequency is a
tolerable proxy, while pinning lets the author assert importance
directly for the specific terms most likely to violate the proxy.

\begin{lstlisting}[caption={Resonance ranking with pinning override}]
for term in vocabulary:
    if term not in draft: continue
    count = occurrences_of(term, corpus)
    if count > 0:
        scored[term] = count

selected = top_k(scored, k = RESONANCE_TOP_K)   # rank-based

for term in pinned_vocabulary:
    if term in draft and term in scored and term not in selected:
        selected.add(term)                      # pin-based override

for term in selected:
    context[term] = surrounding_lines(term, corpus, limit = 10)
\end{lstlisting}

A second, independent source of context is consulted at the same
stage: a persistent ledger file, external to any single run, that
accumulates one row per completed chapter across the pipeline's entire
history.

\begin{lstlisting}[caption={Ledger schema (tab-separated)}]
slug    title    distinction    continuation    outcome    timestamp
\end{lstlisting}

This directly targets the externalize-memory-into-structure pattern:
rather than relying on the corpus directory's raw prose to carry
continuity between chapters, or relying on the author to remember what
an earlier chapter established, the pipeline maintains its own
structured record of what every prior run concluded, and every
corpus-scan stage reads it unconditionally, prepending a summary of
prior chapters' distinctions to the context handed to the critique
stage:

\begin{lstlisting}[caption={Ledger excerpt as injected into critique context}]
### Ledger -- prior chapters on record (externalized memory)
- Repair vs Erasure -- distinction: repair vs erasure (ROOT, CONVERGED)
- Gyrofoil Measure -- distinction: ... (CONTINUES: repair vs erasure, ...)
\end{lstlisting}

The ledger is deliberately flat and append-only rather than queryable
or structured beyond six columns. This is a level-two design choice in
the sense of Section~\ref{sec:levels}: the claim is that externalizing
memory into any inspectable structure, however simple, is better than
not externalizing it at all, not that this particular six-column schema
is the correct one. A richer ledger --- one that also recorded, say,
every definition introduced per chapter rather than only the chapter's
overall distinction --- would strengthen the mechanism further; the
current schema was kept minimal deliberately, to get a working version
of the pattern in place before elaborating it.

\section{Persistence, Convergence, Stalling, and the Invariant Check}
\label{sec:persist}

The critique-revision loop is bounded by a maximum number of
iterations, but is not designed to run for that many iterations
unconditionally. Each critique pass reviews the current draft against
seven checks: unsupported or overstated claims; structural problems
such as weak transitions or redundant sections; hand-wavy reasoning
presented as if it were rigorous; missing definitions for terms used
before they are defined; terminology drift against the vocabulary and
corpus excerpts assembled in Section~\ref{sec:corpus}; theorem or
proposition statements that do not earn their status, such as an
unproven conjecture stated as a theorem; and, targeting the
invariant-search pattern directly, whether the draft explicitly names
something that survives the transformation, reconstruction, or
correction it discusses, or only asserts claims without ever asking
what is preserved. A draft can pass every other check and still fail
this seventh one, if it never actually engages the question the
Distinguish stage raised as its \texttt{INVARIANT\_QUESTION} in the
first place.

Each critique pass is required to end with three lines: an integer
severity score from zero to ten, a qualitative verdict, and a
comma-separated list of which of the seven checks were flagged that
round. The pipeline tracks the severity score across rounds. Two
conditions cause the loop to stop early. If the severity reaches zero,
or the critique explicitly states that no substantive issues remain,
the loop has converged and the current draft is accepted as final. If
the severity instead fails to improve for a configurable number of
consecutive rounds, the loop is judged to have stalled and stops as
well, but with a different message: the last draft before stalling is
reported as needing a human read, rather than being silently accepted
as though the lack of further critique meant the draft had become
adequate.

\begin{definition}
A critique-revision loop is said to \emph{dissolve} at round $n$ if the
severity score has been non-decreasing for $k$ consecutive rounds ending
at $n$, where $k$ is a configured patience parameter, and the loop has
not separately converged. Dissolution is a distinct outcome from
convergence and is reported as such.
\end{definition}

The severity score deserves an explicit caveat. It is produced by
asking the critiquing model to self-report an integer, and nothing in
the pipeline calibrates that integer against any external standard. A
drop from eight to six carries no guarantee of being comparable, in the
amount of actual improvement it represents, to a drop from four to two;
the same numeric gap could correspond to very different amounts of real
progress depending on where in the scale it occurs, or on how the
critiquing model happens to be distributing its scores that round. The
score is a heuristic control variable, used only to distinguish a
downward trend from a flat one, and the pipeline is deliberately built
to ask only that coarse question of it. If the model omits the required
severity line entirely, the score is not left undefined or treated as
zero; it defaults to five, a value chosen specifically so that a
missing score can neither register as convergence nor as a repeated
plateau on its own, forcing the loop to keep going rather than exit on
a signal it never actually received.

\begin{lstlisting}[caption={The persistence loop as a state machine}]
prev = infinity; stall = 0
for round in 1..MAX_ITERATIONS:
    critique = critique(draft)
    sev = parse_severity(critique) or 5     # fallback: never trust silence
    if sev == 0 or says_no_issues(critique):
        return CONVERGED, draft
    stall = stall + 1 if sev >= prev else 0
    prev = sev
    if stall >= DISSOLVE_PATIENCE:
        return DISSOLVED, draft             # last draft before stall, flagged
    draft = revise(draft, critique)
return EXHAUSTED, draft                     # hit MAX_ITERATIONS without either
\end{lstlisting}

Three distinct exit states follow from this, not two. A loop that
converges has a critique attesting the draft is done. A loop that
dissolves has a critique attesting the draft is stuck, with the
specific unresolved issues still on record in the last critique file. A
loop that merely exhausts its iteration budget without doing either ---
possible if severity keeps oscillating rather than trending flat or
down, so neither condition ever triggers --- has neither attestation,
and this is reported as a third kind of ambiguity, deliberately
distinguished from convergence rather than folded into it. An
unbounded loop that merely runs until it exhausts its iteration budget
would conflate several very different situations under a single
stopping condition: a draft that reaches the final round because it
converged three rounds earlier and the loop simply had nothing further
to do is in a different state than a draft that reaches the final
round while still being told the same handful of unresolved problems it
was told rounds earlier, and both are different again from a draft that
oscillated the whole way through. Reporting all three as ``finished''
discards information the author needs; the stall detector and the
exhaustion state both exist to preserve that information rather than
let a fixed iteration count paper over it.

Every critique round's severity and category tags are also appended to
a persistent, cross-run critique log, distinct from the per-run
critique files already written to disk:

\begin{lstlisting}[caption={Critique log schema (tab-separated)}]
slug    round    severity    categories    timestamp
\end{lstlisting}

The purpose of this log, and what it enables beyond this single run, is
taken up in Section~\ref{sec:architecture}.

\section{Turning Criticism into Architecture}
\label{sec:architecture}

The pattern targeted here is specific: rather than treating a piece of
criticism as a local fix, the author's typical response is to promote
it into a mechanism that makes the entire failure mode visible going
forward. A chapter drifting from its intent does not get corrected
once; it produces a fidelity audit that runs on every chapter
thereafter. A critique loop that plateaus without improving does not
get manually interrupted once; it produces a general dissolution
detector. Both of those promotions are already built into this pipeline
as permanent stages, described in Sections~\ref{sec:persist} and
\ref{sec:audit}.

The critique log introduced in Section~\ref{sec:persist} is a
first-order implementation of the mechanism that would let this
promotion happen again in the future, though it is not yet the
promotion itself. Recording slug, severity, and category tags for every
critique round, across every chapter, across the pipeline's entire
history, is what would let a later analysis ask a question no single
run can ask of itself: does the same category of issue --- terminology
drift, missing invariant, illegitimate theorem status --- recur at high
severity across many chapters, rather than being a one-off problem with
one draft. If it does, that recurrence is itself evidence that the
corresponding check deserves to be promoted from a critique-time flag
into an earlier, structural gate, the same way fidelity auditing and
dissolution detection were themselves promoted from ad hoc fixes into
permanent stages of this pipeline.

This essay stops short of specifying that promotion mechanism itself;
the log exists, but nothing yet reads it and proposes a new gate
automatically. This is a deliberate boundary rather than an oversight.
Deciding that a recurring critique warrants a new permanent check is
itself a judgment about what the corpus's failure modes actually are
--- and per Section~\ref{sec:not-do}, judgments of that kind belong on
the human side of this system's boundary, not the automated side. What
the log provides is the raw material for that judgment to be made well
informed, rather than from memory of a handful of recent chapters.

\section{Witnessing the Abstract Structure}
\label{sec:witness}

The asymmetry this stage targets is specific: most writing starts from
examples and generalizes from them; this author's essays typically
start from an abstract structure and only then spend effort finding or
constructing examples that instantiate it, with the examples
functioning as witnesses of an already-established structure rather
than as the motivating cases that produced it.

The Witness stage is a separate model call, run only after the
critique-revision loop has finished, on whatever draft that loop
produced, regardless of which of its three outcomes occurred. It is
given the finished draft and instructed to treat every definition,
proposition, theorem, and lemma as something requiring a witness:
either an existing passage in the draft that already instantiates it
concretely, or, if none exists, one newly constructed concrete example
that does.

\begin{lstlisting}[caption={Witness stage output format}]
WITNESS: Definition 1 -- a torn page reconstructed from ink-transfer
residue witnesses repair preserving structure erasure would destroy.
\end{lstlisting}

This is kept as a separate, late-running pass rather than folded into
drafting for a specific reason: folding it into the draft stage would
reproduce exactly the example-first pattern this stage exists to avoid,
by giving the model the opportunity to reach for a convenient example
before the abstract structure it is meant to witness has actually
stabilized through the critique-revision loop. Running it last ensures
every example is checked against a structure that has already survived
adversarial critique, not one still in flux.

\section{Auditing for Fidelity to the Distinction}
\label{sec:audit}

The final stage compares the finished draft against the original seed
distinction, independent of anything the critique-revision loop said,
and asks specifically whether the finished draft's Crystallized Claim
is still a serious exploration of that same distinction, or whether it
has quietly substituted a different distinction, dropped the original
contrast, or hedged it into vagueness over the course of revision. This
check exists because the critique loop, by construction, only ever
compares a draft to the critique of that draft; nothing earlier in the
pipeline ever looks back at the seed once outlining has finished.
Revision is supposed to sharpen an exploration, not replace it, and a
loop that only measures forward progress against its own critiques has
no way of noticing if it has, cycle by cycle, argued its way to a
different point than the one the author started with. The audit stage
is the only point in the pipeline where the original distinction is
consulted a second time, and it runs against whichever draft the
persistence loop leaves behind regardless of which of the three
outcomes described in Section~\ref{sec:persist} produced it --- a
converged draft can still have drifted, and a dissolved or exhausted
draft is audited too, since a human reading it afterward benefits from
knowing whether it drifted in addition to knowing it stalled.

What counts as drift here is narrower than it might first appear, and
this matters because the draft was never asked to commit to a thesis in
advance. Arriving at a claim no one, including the author, could have
predicted from the seed distinction is expected and correct behavior
for an exploratory draft, not evidence that anything went wrong. What
would count as drift is more specific: quietly substituting a different
distinction than the one the chapter started with, dropping the
original contrast entirely rather than resolving or narrowing it, or
hedging it into vagueness so that the final claim no longer engages the
distinction at all. The audit prompt is written to make this
distinction explicit to the model reviewing it, since a naive
comparison between seed and final draft would otherwise flag every
successful exploration as suspicious simply for not being predictable
in advance.

\section{Relation to the Phoenix Protocol}
\label{sec:phoenix}

Four of the nine stages above --- the admissibility gate, the ranked
retrieval step, the persistence loop, and the fidelity audit --- are
named after, and structurally derived from, the lifecycle of the
Phoenix Protocol: a wave-interference memory architecture, developed as
part of a separate and unrelated project, in which a signal is
admitted into memory subject to an accessibility threshold, persists
across a fixed heartbeat rhythm while its amplitude decays, is
retrieved through a resonance calculation between a query frequency and
the signal's stored frequency, and is finally checked for
reconstruction fidelity against its original form. That architecture is
itself the subject of an extended separate audit, one which treats it
as a speculative computational system to be reconstructed
sympathetically, critiqued rigorously, and selectively reinterpreted in
terms compatible with the author's own framework, rather than accepted
or dismissed wholesale.

The borrowing here is of that same kind, applied to a much smaller
object. Nothing in this pipeline depends on wave mechanics, on a fixed
heartbeat frequency, or on amplitude as a physical or even metaphorical
quantity. What survived the translation was the shape of four
operations --- gate, rank-and-retrieve, persist-with-a-stopping-rule,
audit --- because each one turned out to name a real and useful control
pattern for an iterative language-model loop, independent of the
physical vocabulary it originally arrived in. Renaming a hit-count-based
retrieval ranking as resonance, or a severity plateau as dissolution,
costs nothing and risks a great deal if the names are mistaken for
claims about how the system works. They are not; they are labels kept
for continuity with the architecture that suggested the shape of the
solution, attached to mechanisms that could equally well have been
named after their literal function.

\begin{principle}
A vocabulary may be borrowed from a speculative or unproven system
without borrowing that system's claims, provided every borrowed term is
cashed out, on arrival, as an operation whose behavior can be stated and
tested without reference to the system it came from.
\end{principle}

The remaining five stages --- Distinguish, the exploratory draft's
Crystallized Claim, Dependencies, Witness, and the critique log ---
carry no Phoenix Protocol vocabulary at all. They were derived directly
from the nine writing patterns described in Section~\ref{sec:levels}'s
motivation, with no intermediate translation step through any other
system. The Phoenix-derived stages and these five are therefore
independent in origin, and their coexistence in one pipeline is a claim
only about what turned out to be jointly useful, not about any deeper
relationship between the two sources.

\section{The Continuation Declaration as Human Decision Surface}
\label{sec:continuation-decl}

The pipeline drafts, checks internal and corpus-level consistency,
checks against a persistent ledger of prior chapters, revises against
critique, checks for a named invariant, detects its own stalling,
discovers and queues missing dependencies, finds witnesses for its own
abstract structure, and checks for drift from original intent. What it
does not do, at any stage, is decide what a chapter is about, which of
several candidate primitives deserves to be treated as foundational,
whether a queued dependency is actually worth writing, whether a given
chapter is a root contribution or a continuation of existing work, or
how two independently developed concepts ought to relate once both
exist. Those decisions constitute what might be called the human
decision surface of the larger research program, and every one of them
either happens before a seed ever reaches the Distinguish stage, or
happens when a human reads a finished chapter, a queued dependency, or
a ledger entry and judges it.

Chapter accretion, described in Section~\ref{sec:accretion}, is the
part of this pipeline most likely to have quietly crossed that
boundary, since a system that automatically discovers and queues new
chapters is one plausible reading away from a system that decides, on
its own, what the corpus needs next. The boundary is preserved by a
specific, narrow mechanism rather than by general good intentions: a
seed file's optional first line may declare \texttt{ROOT} or
\texttt{CONTINUES: <terms>}, and this declaration is read, never
inferred. The dependency stage populates the \texttt{CONTINUES:} line
of every queued seed with the chapter that spawned it, so provenance is
never lost, but it does not decide whether that queued concept is
actually foundational, actually worth its own chapter, or actually
correctly scoped --- those questions remain open until a human either
runs the queued seed through the pipeline as written, edits it first,
or discards it without running it at all. A run with no continuation
declaration at all proceeds --- the pipeline does not block on it ---
but is logged as \texttt{UNSPECIFIED} rather than having the gap
silently filled by a guess.

\begin{principle}
A system that discovers structure automatically (missing dependencies,
recurring critique categories, resonant vocabulary) may surface that
structure as a proposal, logged and inspectable, without thereby being
entitled to decide the proposal is correct. The line between
discovering and deciding is the line this system is required to keep
on the human side of, at every stage where automation is added.
\end{principle}

\section{What the Pipeline Does Not Do}
\label{sec:not-do}

It is worth stating the resulting boundary plainly, in one place, since
it is easy to lose track of across nine stages. The pipeline does not
decide what a chapter is about, which of several candidate primitives
deserves to be treated as foundational, whether a queued dependency is
actually worth writing, whether a chapter is a root or a continuation
of existing work, how two independently developed concepts ought to
relate once both exist, or whether a recurring critique category
warrants promotion into a new permanent gate. It does not notice
scoping errors. A critique-revision loop can be run indefinitely over a
chapter built on a wrongly chosen foundational distinction and it will
only ever produce a more polished chapter built on the wrong
foundation; accretion can run for many recursive levels against a
wrongly scoped root distinction and will only ever produce a larger
structure built on that wrong scoping. Nothing in this specification is
capable of noticing that the foundation or the scoping was wrong,
because noticing that requires comparing the chapter, or the emerging
graph of chapters, against the rest of the program's architecture ---
precisely the part this system was not built to hold. What it was built
to hold is narrower and, within that scope, thorough: drafting,
internal and cross-corpus consistency, convergence detection, missing
dependency discovery, concrete instantiation of abstract structure, and
fidelity to stated intent, all logged, none of it silently discarded,
and all of it subordinate to decisions a human makes before, during,
and after every run.

\section{Conclusion: Continuation, Not Generation}
\label{sec:conclusion}

Read back across the preceding sections, a single principle organizes
what would otherwise be a list of unrelated engineering choices. The
distinction-extraction stage of Section~\ref{sec:distinguish} exists so
that a chapter is answerable to something the author already judged
worth preserving, rather than to a thesis invented for the occasion.
The vocabulary constraint of Section~\ref{sec:vocab} exists so that new
writing is checked against what earlier writing already meant. The
ledger of Section~\ref{sec:corpus} exists so that a corpus-scan reads
the pipeline's own accumulated history rather than starting from a
clean slate on every run. The dependency-discovery mechanism of
Section~\ref{sec:accretion} exists so that a missing foundation is
queued to be actually built rather than patched over in the chapter
that needed it. The invariant check of Section~\ref{sec:persist} exists
so that no chapter is allowed to claim anything without having said,
somewhere, what survives the transformation it concerns. The audit of
Section~\ref{sec:audit} exists so that the finished chapter can be
checked against the intent that started it, not only against its own
intermediate critiques. The continuation declaration of
Section~\ref{sec:continuation-decl} exists so that the pipeline's own
growth is honest about whether it is starting something new or
extending something already underway. In every case, something already
committed to functions as a constraint on what may legitimately come
next, and none of the nine stages generates content from an
unconstrained space.

This is a level-three claim in the sense of Section~\ref{sec:levels},
and it should be read as a hypothesis about this pipeline and about
programs structured like it, not as an established fact about writing
systems in general. But it is a hypothesis that follows directly, at
the level of mechanism rather than metaphor, from commitments that
organize the author's broader theoretical program: history before
state, repair before correctness, admissibility before prediction.
Those commitments say, at the level of a research program, that what
has already happened constrains what is admissible next, that fixing
an error takes priority over asserting a clean answer, and that whether
a continuation is allowed at all takes priority over predicting which
continuation is likeliest. This pipeline is the same claim applied one
level down, to the mechanics of producing a single chapter and, through
accretion, a small graph of them: what has already been written
constrains what a draft may claim, fixing a draft's problems takes
priority over letting it stand, and whether a draft is even admissible
for revision takes priority over predicting how good a revision of it
might be. The pipeline does not implement history-before-state,
repair-before-correctness, or admissibility-before-prediction as formal
objects; it implements the much smaller, purely mechanical shadow those
commitments cast on the specific problem of drafting prose, with a
language model in the loop, inside a corpus that is still growing.

One genuine bug in the mechanism meant to let these nine stages compose
was found by testing during the development of this pipeline, rather
than assumed away, and is documented in Section~\ref{sec:incident}
rather than quietly patched: an identity function safe for a single
chapter collapsed once dependency chains composed it with itself
repeatedly, and the fix required asking, of every identity the system
constructs, whether it survives the specific transformation --- being
prefixed onto again, recursively --- that this system actually performs.
That the bug was found this way, in the pipeline's own source rather
than in someone else's prose, is itself a small instance of the larger
claim: a system built to ask what survives a transformation eventually
has to ask that question of itself, and answering it honestly, in
public, in the specification rather than only in a private fix, is what
keeps continuation as a real commitment rather than a stated one. The
pipeline continues a corpus; it does not decide what the corpus is
about.

\end{document}

